% [original page: 32]
\noindent מהמשפטים האחרונים נובע מיד

\bigskip
\noindent\textbf{\underline{משפט 5.}}

\noindent התנאי ההכרחי והמספיק לכך שמספר $m$ מסויים יותן להצגה
עי"י התבנית $(a,b,c)$ באופן עיקרי, בתנאים האמורים, הוא
שהמספר $1$ יותן להצגה עי"י התבנית
\[ \left(am^{2\nu-1},\; Zm^{\nu-1},\; n\right) \]
בעלת המבחין $m^{2(\nu-1)}D$,

\noindent במקום שהמספרים $Z, n$ נקבעים עפ"י הנוסחאות (13)(14)(17):

\noindent מספר טבעי כל שהוא.

% [original page: 33]
\subsection*{3\S. תכונות הפתרונות בקבוצה אחת.}

\noindent 1. כפי שיתברר להלן נשמרות התכונות, שהזכרנו בפרק ב' בטיטת
הקלאספיקציה של גאוס, ועוד אחרות, בתבניות ריבועיות הן לגבי
מייצג $m$ אחד נתון והן לגבי שני מייצגים $m, m'$ שונים בתנאים
מסויימים. על החלוקה לקבוצות של פתרונות של תבנית התבנית, בנות
התכונות הנ"ל נעמוד בסעיף זה.

\noindent נוכיח באמת וראשונה את המשפט הבא:

\bigskip
\noindent\textbf{\underline{משפט 6.}}

\noindent בשביל כל זוג פתרונות $(x,y), (x',y')$ של התבנית
$(a,b,c)$, הן ביחס לאותו מייצג $m$ והן ביחס לשני
מייצגים $m, m'$ שונים בתנאים:
\[ (1) \qquad m = m_1 m_2 \qquad\qquad |m_1| > 1 \]
\[ m' = m_1 m_2' \]

\noindent הסמוכים לטרסים $r, er_1$ $(e = \pm 1)$ של הקונגרואנציה.
\[ (2) \qquad ar^2 + br + c \equiv 0 \pmod{m_1} \]

\noindent מתקיים:
\[ (3) \qquad xy' - ex'y \equiv yy'(r - er_1) \pmod{m_1} \]

\noindent ולהיפך: כל זוג פתרונות $(x,y), (x',y')$ שמקיימות את התנאי
האחרון שייכות לקבוצות הפתרונות המיוחסות לטרסים $r,$
$er_1$ $(e=\pm 1)$ הנ"ל.

\bigskip
\noindent\textbf{\underline{הוכחה}}

\noindent לפי הכתוב נוכל ליצור, כפי שהתברר בסעיפים הקודמים את
הקטרים:
\[ (4) \qquad x = m_1 q + ry \]
\[ x = m_1 q' + r_1 y' \]

\noindent בתנאים:
\[ 0 \le r < |m_1| \]
\[ 0 \le r_1 < |m_1| \]

\noindent אולם מכאן נקבל מיד
\[ xy' - ex'y = (m_1 q + ry)y' - {} \]

% [original page: 34]
\[ {} - e(m_1 q' + r_1 y')y = m_1(qy' - eq'y) + yy'(r - er_1) \]

\noindent ולכן נופק מיד, מ.ט.כ.:
\[ \text{מ.ט.ל.} \qquad xy' - ex'y \equiv yy'(r - er_1) \pmod{m_1} \]

\bigskip
\noindent\textbf{\underline{מסקנה 1.}} כאשר
\[ \left. \begin{array}{l} (yy', m) = 1 \Longleftrightarrow (yy', m_1) = 1 \\ r - er_1 \equiv 0 \pmod{m_1} \end{array} \right\} \]

\noindent אזי
\[ (5) \qquad xy' - ex'y \equiv 0 \pmod{m_1} \]

\bigskip
\noindent\textbf{\underline{מסקנה 2.}} כאשר
\[ e = 1 \, , \qquad m' = m_1 = m \, , \qquad \nu = 1 \]

\noindent אנו מקבלים
\[ xy' - x'y \equiv 0 \pmod{m} \]

\noindent כלומר את המסקנה הידועה מהקלאספיקציה של גאוס (פרק ב').

\medskip
\noindent 2. נטפל עתה במקרה
\[ m = m' = m_1 \, . \]

\noindent לפי ההנחה
\[ (6) \qquad x = m q + ry \]

\noindent אנו מקבלים מיד, בדרך שהבאנו בסעיף 2, את המשוואה
\[ (7) \qquad am^{2\nu-1}q^2 + Zm^{\nu-1}qy + ny^2 = 1 \]

\noindent לפי המסקנות שנתקבלו בסעיף הקודם אפשר, איפוא, לקבל את כל
הפתרונות $(x,y)$ של התבנית $(a,b,c)$ ביחס למיוצג $m$
המסויכים לטורס $r$ של הקונגרואנציה
\[ (8) \qquad ar^2 + br + c \equiv 0 \pmod{m} \]

\noindent בתנאי
\[ 0 \le r < |m| \]

\noindent מהמשוואה (7) באמצעות הקשר (6).

\noindent אני טוען שהתבנית
\[ \left(am^{2\nu-1},\; Zm^{\nu-1},\; n\right) \]

\noindent אקוויוולנטית לתבנית אחרת $m$ מהצורה

% [original page: 35]
\[ (9) \qquad (1,\; Z',\; n') \, , \]

\noindent כשהמקדמים $Z', n'$ מקיימים את הנוסחאות:
\[ (10) \quad \left\{ \begin{array}{l}
Z' = 2am^{2\nu-1}q_1\beta + Zm^{\nu-1}(q_1\delta + y_1\beta) + 2ny_1\delta \\[4pt]
n' = am^{2\nu-1}\beta^2 + Zm^{\nu-1}\beta\delta + n\delta^2 \, ,
\end{array} \right. \]

\noindent במקום ש-$(q_1, y_1)$ הנו פתרון מסויים של המשוואה (7), ומתקיים
בבת-אחת:
\[ (11) \qquad \delta q_1 - \beta y_1 = 1 \]

\noindent ואת, כל הפתרונות $(q,y)$ של המשוואה (7) אפשר לקבל מהנוסחאות
\[ (12) \quad \left\{ \begin{array}{l}
q = q_1 \bar{q} + \beta w \\
y = y_1 \bar{q} + \delta w
\end{array} \right. \]

\noindent במקום ש-$\bar{q}, w$ הם הפתרונות של המשוואה
\[ (13) \qquad \bar{q}^2 + Z'\bar{q}w + n'w^2 = 1 \, . \]

\noindent כל המסקנות הללו נובעות במישרים עפ"י שיטתו של גאוס, שהובאה
בפרק ב'. מכאן, שבכדי למצוא את כל קבוצת הפתרונות של תבנית
$(a,b,c)$ מסויימת ביחס למיוצג $m$ מסויים עלינו לדעת את
אחד הפתרונות של המשוואה (7) ולקבוע שיטה לפתירת המשוואה (13)
האחרונה. בדבר הראשון נעסוק בפרק הבא, בהנחה, כי:
\[ D < 0 \]

\noindent ובדבר הטכני נדון במסגרת הענין כאן.

\medskip
\noindent 3. לשם פתירת המשוואה (13) נניח:
\[ (14) \quad \left\{ \begin{array}{l}
2\bar{q} + Z'w = u \\
w = v
\end{array} \right. \]

\noindent ברור שהמספרים $u, v$ ישארו שלמים עפ"י הנחה זאת. ומאידך נקבל
עי"י הצגת הערכים
\[ (15) \quad \left\{ \begin{array}{l}
w = v \\[4pt]
\bar{q} = \dfrac{u - Z'v}{2}
\end{array} \right. \]

\noindent אל התבנית (13) נקבל מיד את משוואת הידועה

% [original page: 36]
\[ (16) \qquad \boxed{\, u^2 - m^{2(\nu-1)}Dv^2 = 4 \,} \]

\noindent פתירת המשוואה (13) מתלכדת, איפוא, עם פתירת משוואת פל: שתי
המשוואות הנ"ל, איפוא, רק שתי צורות שונות של משוואה אחת
והיא משוואת פל.

\bigskip
\noindent\textbf{\underline{הערה.}} הסימון היחיד שנתקבל כאן במשוואת פל הוא בזה שבמקום $D$
מופיע $m^{2(\nu-1)}D$ אולם כפי שמתברר בתורת משוואת פל יש תמיד
פתרון למשוואה
\[ u^2 - m^{2(\nu-1)}Dv^2 = 4 \qquad \text{(בשביל כל } m, \nu \text{)} \]

\noindent כאשר יש פתרון למשוואה
\[ u^2 - Dv^2 = 4 \, . \]

\bigskip
\hrule
\bigskip

% [original page: 37]
\section*{פרק ד'.}

\subsection*{1\S. פתירת תבנית היחידה כאשר $D < 0$}

\noindent 1. כפי שראינו מצטמצמת בעיית פתירתה של תבנית ריבועית $(a,b,c)$
כל שהוא לגבי מייצג $m$ מסויים לידי פתרות תבנית ריבועית מסויימת
לגבי המייצג $1$.

\noindent כל תבנית ריבועית המייצגת את המספר $1$ נקראת בשם תבנית היחידה
והבעייה העומדת לפנינו עתה, היא: פתירתה של תבנית היחידה.

\noindent בפרק זה נטפל רק במקרה המיוחד כאשר
\[ (1) \qquad D < 0 \]

\noindent יהי, איפוא, נתון:
\[ (2) \qquad ax^2 + bxy + cy^2 = 1 \]

\noindent בתנאי (1).

\noindent מעצם קיומו של השוויון (2) במספרים שלמים $(x,y)$ נובע:
\[ (3) \qquad (x,y) = 1 \]

\noindent אולם, לפי"ז אפשר למצוא שני מספרים $\beta, \delta$ שלמים המקיימים את-השוויון:
\[ (4) \qquad \delta x - \beta y = 1 \]

\noindent נשתמש עתה בתחליף מהצורה:
\[ (5) \quad \left\{ \begin{array}{l}
x = x_1 s + \beta t \\
y = y_1 s + \delta t \, ;
\end{array} \right. \]

\noindent במקום ש-$(x_1, y_1)$ הוא פתרון מסויים של השוויון (2).

\noindent עי"י תחליף זה עובר השוויון (2) לשוויון:
\[ (6) \qquad s^2 + zst + nt^2 = 1 \]

\noindent במקום ש-$Z, n$ מקיימים את התנאים:
\[ (7) \quad \left\{ \begin{array}{l}
z^2 \equiv D \pmod{4} \\[4pt]
n = \dfrac{z^2 - D}{4}
\end{array} \right. \]

\noindent כפי שנובע מן ההנחה

% [original page: 38]
\[ (8) \qquad z_k = (2ax+by)\beta_k + (bx+2cy)\delta_k \]

\noindent במקום :-
\[ (9) \quad \left\{ \begin{array}{l}
\beta_k = \beta_0 + kx \\
\delta_k = \delta_0 + ky
\end{array} \right. \]

\noindent ו-$(\beta_0, \delta_0)$ הנו פתרון מסויים של המשוואה (4), בשביל כל $k$ שלם.

\noindent כי מתוך הנוסחאות הנ"ל אנו מקבלים:
\[ z_k = z_0 + (2ax+by)kx + (bx+2cy)ky = z_0 + 2k \]

\noindent לפי (2).

\noindent מאידך אנו למדים מן השוויון האחרון, שבין כל השרשים $z$ קיים
אחד ויחיד המקיימים את התנאים:
\[ (10) \qquad 0 \le z < 2 \, . \]

\noindent אולם
\[ D = b^2 - 4ac \equiv 0, 1 \pmod{4} \]

\noindent בהתאם לכך, אם $b$ הנו זוגי או לא זוגי.

\noindent מכאן, לקונגרואנציה
\[ z^2 \equiv D \pmod{4} \]

\noindent ישנו תמיד שורש אחד ויחיד ברווח (10), והוא $0$ או $1$ בהתאם לכך אם
$b$ הנו זוגי או לא-זוגי.

\noindent אנו מקבלים, איפוא, את המשפט הבא:

\bigskip
\noindent\textbf{\underline{משפט 1.}}

\noindent תבנית היחידה היא אקוויוולנטית לתבנית
\[ (11) \qquad \left(1,\, 0,\, -\dfrac{D}{4}\right) \]

\noindent או לתבנית
\[ \left(1,\, 1,\, \dfrac{1-D}{4}\right) \]

\noindent בהתאם לכך אם $b$ הנו זוגי או לא זוגי.

\noindent ואמנם מתוך התחליף (5) בהנחה (4) נובע, כי:
\[ z^2 - 4n = D \]

% [original page: 39]
\noindent ולכן
\[ n = \frac{0 - D}{4} \]

\noindent או
\[ n = \frac{1 - D}{4} \]

\noindent כנ"ל.

\medskip
\noindent 2. הבעייה העיקרית של פתירת תבנית היחידה לגבי $1$ מצטמצמת, מעתה,
לידי מציאת מקדמי הטרנספורמציה (5): כי מבין אלו המקדמים העריכים
$(x_1, y_1)$ הם המקיימים את ההצגה (2).

\noindent נכתוב את השוויון (2) בצורה
\[ (12) \qquad (2ax+by)^2 - Dy^2 = 4a \]

\noindent ונכיח, $D<0$, כי:

\noindent מן השוויון האחרון אנו מסיקים מיד, לפי ההנחה שעשינו, כי:
\[ (13) \quad \left\{ \begin{array}{l}
4a > 0 \\
a > 0 \, .
\end{array} \right. \]

\noindent ולכן כי:

\noindent אולם את השוויון (12) אפשר גם לכתוב בצורה
\[ (2cy+bx)^2 - Dx^2 = 4c \]

\noindent ולכן נוכל גם להסיק, כי:
\[ (14) \qquad c > 0 \, . \]

\noindent נשתמש עתה בתחליף
\[ (15) \quad \left\{ \begin{array}{l}
x = \alpha_1 x' + \beta_1 y' \\
y = \gamma_1 x' + \delta_1 y'
\end{array} \right. \]

\noindent בהנחה, כי:
\[ (16) \qquad \alpha_1 = 0 \, , \quad \beta_1 = 1 \, , \quad \gamma_1 = -1 \]

\noindent ברור, שלפי הנחה זאת מתקיים
\[ (17) \qquad \alpha_1 \delta_1 - \gamma_1 \beta_1 = 1 \]

\noindent ולכן תעבור עי"י תחליף זה התבנית (2) לתבנית $(a_1, b_1, c_1)$
אחרת האקוויוולנטית לה, והמקדמים החדשים יתקבלו ממערכת השוויונים:

% [original page: 40]
\[ (18) \quad \left\{ \begin{array}{l}
a_1 = c \\
b_1 = -b - 2c\delta_1 \, ,
\end{array} \right. \]

\noindent כפי שקל להסיק מחשבונות רגילים. עד כה לא עשינו כל הגבלות לגבי
$\delta_1$, ולכן בחירתו היא חפשית. נקבע את ערכו באופן שיתקיים:
\[ (19) \qquad -2a_1 < b_1 < 2a_1 \, . \]

\noindent לשם כך מספיק לבחור בו באופן שיתקיים
\[ (20) \qquad -\frac{b}{2c} - 1 < \delta_1 < 1 - \frac{b}{2c} \]

\noindent כלומר עלינו לקבוע את $\delta_1$ עפ"י הנוסחה
\[ (21) \qquad \boxed{\, \delta_1 = -E\!\left(\dfrac{b}{2c}\right) \,} \]

\noindent במקום ש-$E(n)$ מסמן את החלק השלם של המספר $n$.
\[ \text{הטרנספורמציה} \qquad \begin{pmatrix} \alpha_1 & , & \beta_1 \\ \gamma_1 & , & \delta_1 \end{pmatrix} \]

\noindent המקיימת את הדרישות הנ"ל היא, איפוא, קבועה באופן חד-ערכי עי"י
המקדמים $b, c$ של התבנית הנתונה, ומכאן:

\bigskip
\noindent\textbf{\underline{משפט 2.}}

\[ \text{עי"י הטרנספורמציה} \qquad \begin{pmatrix} 0 & , & 1 \\ -1 & , & -E\!\left(\dfrac{b}{2c}\right) \end{pmatrix} \]

\noindent עוברת התבנית $(a,b,c)$ המייצגת את המספר $1$ לתבנית אחרת
$(a_1, b_1, c_1)$ אקוויוולנטית לה, שבשביל מקדמיה מתקיימים התנאים:
\[ -2a_1 < b_1 < 2a_1 \, . \]

\noindent אולם אפשר להמשיך ולהשתמש בטרנספורמציות מהסוג הקודם גם הלאה.
לפיכך נקבל תבנית $(a_2, b_2, c_2)$ חדשה בתנאים:
\[ \left\{ \begin{array}{l}
a_2 = c_1 \\
-2a_2 < b_2 < 2a_2
\end{array} \right. \]

% [original page: 41]
\noindent ובאופן כללי, נקבל מערכת של תבניות אקוויוולנטיות
\[ \left\{ \begin{array}{l}
a_\nu = c_{\nu-1} \\
-2a_\nu < b_\nu < 2a_\nu
\end{array} \right. \qquad \text{בשביל} \qquad \nu = 3, 4, \ldots \]

\noindent אני טוען שמוכרח להתקיים אחד התנאים
\[ (22) \quad \left\{ \begin{array}{l}
a_\nu \le c_{\nu-1} \\
a_\nu = c_\nu
\end{array} \right. \qquad \text{או} \]

\noindent ואמנם, לו היה
\[ a_\nu > c_\nu \]

\noindent היה יוצא, כי:
\[ a_{\nu+1} = c_\nu < a_\nu \qquad \text{בשביל כל} \qquad \nu = 3, 4, \ldots \]

\noindent ולכן היינו מקבלים סדרה אין סופית של מספרים טבעיים בתנאים
\[ a_1 > a_2 > a_3 > \cdots \]

\noindent מאחר, שכפי שהוכחנו כבר קיים בשביל המקדמים הקיצוניים של התבנית
\[ a_\nu > 0 \, , \qquad c_\nu > 0 \qquad \text{בשביל כל} \qquad \nu = 1, 2, 3, \ldots \]

\noindent ודבר זה מתנגד למציאות איבר מינימלי בכל סדרה של מספרים טבעיים.
מכאן -- נכונות הטענה, כלומר קיום אחד התנאים (22), בשביל $\nu$ מסויים.

\medskip
\noindent 3. נניח עתה שקיים הראשון מתנאי (22). עי"י התהליך המתבסס על סורת
טרנספורמציות מהסוג הנקבע במשפט 2 אנו מגיעים לבסוף לתבנית מהצורה
\[ (23) \qquad a_\nu x^{(\nu)2} + b_\nu x^{(\nu)} y^{(\nu)} + c_\nu y^{(\nu)2} = 1 \]

\noindent בתנאים
\[ (24) \quad \left. \begin{array}{l}
|b_\nu| < 2a_\nu \\
a_\nu \le c_\nu - 1
\end{array} \right\} \]

\noindent מתנאים אלו:
\[ b_\nu^2 < 4a_\nu^2 \]
\[ b_\nu^2 < 4a_\nu(c_\nu - 1) \]

% [original page: 42]
\noindent ולכן
\[ b_\nu^2 - 4a_\nu c_\nu < -4a_\nu \]

\noindent או
\[ D < -4a_\nu \]

\noindent ולכן
\[ (25) \qquad -D > 4a_\nu \, . \]

\noindent נכתוב את השוויון (23) בצורה:
\[ (2a_\nu x^{(\nu)} + b_\nu y^{(\nu)})^2 - Dy^{(\nu)2} = 4a_\nu \]
\[ Dy^{(\nu)2} = (2a_\nu x^{(\nu)} + b_\nu y^{(\nu)})^2 - 4a_\nu \]

\noindent מכאן כי:
\[ y^{(\nu)2} \le -\frac{4a_\nu}{D} \qquad \left( \text{כאשר } a_\nu > 0 \, , \; D < 0 \right) \]

\noindent ולכן נסיק מיד על סמך ההנחה (25), כי:
\[ (26) \qquad \boxed{\, y^{(\nu)} = 0 \,} \]

\noindent אולם, לפי"ז נובע מהשוויון (23), כי:
\[ (27) \qquad \boxed{\, a_\nu = 1 \,} \]

\noindent ולפיכך בהתאם ל-(24):
\[ (28) \quad \left\{ \begin{array}{l}
b_\nu = 0 \\
b_\nu = \pm 1
\end{array} \right. \qquad \text{או} \]

\noindent בהתאם לכך, אם
\[ \left\{ \begin{array}{l}
D \equiv 0 \pmod{4} \\
D \equiv 1 \pmod{4} \, .
\end{array} \right. \qquad \text{או} \]

\noindent הגענו, איפוא, לידי מסקנה, שאם קיים הראשון מתנאי (22) אזי
צורת הטרנספורמציות המבוצעות עפ"י משפט 2, הן, הן הקובעות את
הטרנספורמציה המעבירה את התבנית $(a,b,c)$ לאחת התבניות
\[ \left(1,\, 0,\, -\dfrac{D}{4}\right) \qquad \text{או} \]
\[ \left(1,\, \pm 1,\, \dfrac{1-D}{4}\right) \]

\noindent בהתאם לכך, אם
\[ D \equiv 0 \pmod{4} \]

% [original page: 43]
\noindent או
\[ D \equiv 1 \pmod{4} \, , \]

\noindent ובדרך זאת קבענו מצד אחד שיטה לביצוע הטרנספורמציה הנזכרת
במשפט 1, ומאידך גיסא פתרנו את המשוואה (23).

\medskip
\noindent 4. לבסוף נניח, שקיים השני מהתנאים (22). המשוואה (23) מקבלת
אז את הצורה:
\[ (29) \qquad a_\nu x^{(\nu)2} + b_\nu x^{(\nu)} y^{(\nu)} + a_\nu y^{(\nu)2} = 1 \]

\noindent אולם ללא כל הגבלה מתקיים
\[ \left(x^{(\nu)} \pm y^{(\nu)}\right)^2 \ge 0 \]

\noindent ולכן גם
\[ (30) \qquad a_\nu x^{(\nu)2} \pm 2a_\nu x^{(\nu)} y^{(\nu)} + a_\nu y^{(\nu)2} \ge 0 \]

\noindent מכיוון ש:
\[ a_\nu > 0 \]

\noindent נראה עתה, שפתרון המשוואה (29) אינו מהווה קושי רציני.
נבדיל לשם כך בין שני מקרים:
\[ \text{(א} \qquad x_\nu y_\nu < 0 \]
\[ \text{(ב} \qquad x_\nu y_\nu > 0 \]

\noindent \textbf{א.} במקרה א' אינו קיים הסימן העליון בנוסחה (30).
מכיוון, שלפי ההנחה
\[ |b_\nu| < 2a_\nu \]

\noindent היה יוצא מיד, כי:
\[ a_\nu x^{(\nu)2} + b_\nu x^{(\nu)} y^{(\nu)} + a_\nu y^{(\nu)2} >
a_\nu x^{(\nu)2} + 2a_\nu x^{(\nu)} y^{(\nu)} + a_\nu y^{(\nu)2} \]

\noindent ולכן
\[ a_\nu x^{(\nu)2} + b_\nu x^{(\nu)} y^{(\nu)} + a_\nu y^{(\nu)2} > 1 \]

\noindent בניגוד ל-(29).

\noindent לפיכך מוכרח להיות
\[ a_\nu \left(x^{(\nu)} + y^{(\nu)}\right)^2 = 0 \]

% [original page: 44]
\noindent ומכאן
\[ (31) \qquad \boxed{\, x^{(\nu)} + y^{(\nu)} = 0 \,} \]

\noindent \textbf{ב.} במקרה ב' נופק, כי:
\[ a_\nu x^{(\nu)2} + b_\nu x^{(\nu)} y^{(\nu)} + a_\nu y^{(\nu)2} >
a_\nu x^{(\nu)2} - 2a_\nu x^{(\nu)} y^{(\nu)} + a_\nu y^{(\nu)2} \]

\noindent ולכן לו היה
\[ a_\nu \left(x^{(\nu)} - y^{(\nu)}\right)^2 > 0 \]

\noindent היינו מקבלים
\[ a_\nu x^{(\nu)2} + b_\nu x^{(\nu)} y^{(\nu)} + a_\nu y^{(\nu)2} > 1 \]

\noindent בניגוד ל-(29).

\noindent לפיכך קיים במקרה זה:
\[ (32) \qquad \boxed{\, x^{(\nu)} - y^{(\nu)} = 0 \,} \]

\noindent הוכחנו, איפוא, שבשני המקרים א', וב', קיים:
\[ (33) \qquad \boxed{\, x^{(\nu)} = \pm y^{(\nu)} \,} \]

\noindent בזה הצלחנו לפתור את השוויון (23) בכל המקרים ולכן על סמך
שיטת הטרנספורמציות הכלולה במשפט 2, הצלחנו גם לפתור את המש-
וואה (2) בכל המקרים.

\medskip
\noindent 5. נטפל עוד עתה בחקירת מבנה התבנית (29) בתנאים האמורים.
עי"י הצבת הערכים (33) לתבנית אנו מקבלים מיד
\[ x^{(\nu)2}(2a_\nu \pm b_\nu) = 1 \]

\noindent ומכאן נופק מיד, כי:
\[ (34) \qquad \boxed{\, 2a_\nu \pm b_\nu = 1 \,} \qquad \boxed{\begin{array}{l} x^{(\nu)} = \pm 1 \\ y^{(\nu)} = \pm 1 \end{array}} \]

\noindent למקרה זה יש, איפוא, מקום, רק כאשר $b_\nu$ הנו בלתי-זוגי; ובכן,

% [original page: 45]
\[ (35) \qquad \boxed{\, D \equiv 1 \pmod{4} \,} \qquad \text{כאשר} \]

\noindent חוץ מזה יש מקום לשוויון (34) רק בשביל הסימן התחתון, הואיל
ו-$a_\nu > 0$ ;

\noindent לפיכך נקבל:
\[ (36) \qquad \boxed{\, a_\nu = \dfrac{b_\nu + 1}{2} \,} \]

\medskip
\noindent 6. לפי משפט 2 (29) התבניות הן אקוויוולנטיות ולכן קיים:
\[ b_\nu^2 - 4a_\nu^2 = D \]

\noindent או, לפי (36)
\[ b_\nu^2 - (b_\nu + 1)^2 = D \]

\noindent כלומר
\[ D = -(2b_\nu + 1) \]

\noindent אולם לפי (35)
\[ D = 4\Delta + 1 \]

\noindent ולכן
\[ 4\Delta = -2b_\nu - 2 \]

\noindent כלומר
\[ (37) \qquad \boxed{\, b_\nu + 1 = -2\Delta \,} \]

\noindent את התבנית (29) אפשר, איפוא, לקבוע מראש, במקרה ב', עפ"י
הנוסחאות:
\[ (38) \quad \left\{ \begin{array}{l}
a_\nu = -\Delta \\
b_\nu = -2\Delta - 1 \\
D = 4\Delta + 1 \, .
\end{array} \right. \qquad \text{כאשר} \]

\bigskip
\hrule
\bigskip

% [original page: 46]
\section*{פרק ה'.}

\subsection*{על קונגרואנציות ממעלה שניה}

\subsection*{1. הערות על קונגרואנציות ממעלה שניה.}

\noindent 1. כפי שראינו בפרקים ב', וג', מובאת בעיית פתירת התבנית הריבועית
לגבי $m$ מסויים, לפי הקלאספיקציה של גאוס לידי פתרון הקונגרואנציה
הלא-מלאה
\[ (1) \qquad z^2 \equiv D \pmod{4m} \]

\noindent או, לפי הקלאספיקציה החדשה לידי פתרון הקונגרואנציה המלאה
\[ (2) \qquad ar^2 + br + c \equiv 0 \pmod{m} \]

\noindent בשביל $m$ מסויים.

\noindent סומה עלינו, איפוא, לברר כמה פרטים, מתורת הקונגרואנציות, הנוגע-
ים במישרים לתורת התבניות.

\medskip
\noindent\textbf{\underline{2. הצגת הבעייה.}}

\noindent בתורת הקונגרואנציות כשלעצמה אנו נתקלים כבר בבעיות היסוד דלקמן:

\noindent (א שיטת פתירת כל אחת הקונגרואנציות (1) (2):

\noindent (ב) איך לקבוע את מספר השרשים הפרימיטיביים של הקונגרואנציות
(1) (2).

\noindent (ג מהו הקשר בין הקונגרואנציות (1) (2)?

\noindent (ד) מתי קיימים לקונגרואנציה (2) שרשים המקיימים את אחד הת-
נאים:
\[ (r_1 - r_2,\, m) = 1 \qquad \text{או} \]
\[ (r_1 + r_2,\, m) = 1 \, . \]

\noindent\textbf{\underline{הערה.}} לבעייה ד' יש ערך רב, ערך שכבר ראינו בקלאספיקציה החדשה
של פתרונות התבנית הריבועית ועוד נשוב לטפל בענין זה באחד הפר-
קים דלהלן.

\medskip
\noindent\textbf{\underline{3. קונגרואנציות לא מלאות.}}

\noindent תורת הקונגרואנציות מהצורה

% [original page: 47]
\[ (3) \qquad z^2 \equiv D \pmod{m} \]

\noindent ידועה בשם ``תורת השאריות הריבועיות''. תורה זאת פותחה ועובדה כהלכה
עי"י גדולי המתימטיקאים, כגון לז'נדר, גאוס, יעקבי, דיריכלה, איזנ-
שטיין ועוד.

\noindent לעצם תורה זאת לא ניכנס כאן ונסתפק רק בקביעת מספר השרשים הי-
סודיים של הקונגרואנציה (3), מאחר שעובדה זאת תהיה חשובה לנו להבא,
לשם הערכת מספר קבוצות הפתרונות השונות אשר לתבנית ריבועית. והנה
המסקנות הכלליות הידועות בשטח זה:

\bigskip
\noindent\textbf{\underline{משפט 1.}} מספר שרשי הקונגרואנציה (3) הוא:
\[ 2^{\lambda+\mu} \]

\noindent במקום ש-$\lambda$ הוא מספר הגורמים הראשוניים השונים של $m$;
ובמקום
\[ \begin{array}{lll}
4 \nmid m & \text{כאשר} & \mu = 0 \\
m = 4(2k+1) & \text{כאשר} & \mu = 1 \\
m = 8k & \text{כאשר} & \mu = 2
\end{array} \]

\noindent בתנאי:
\[ (D, m) = 1 \]

\noindent למסקנות הכלולות במשפט זה אנו מגיעים על סמך כמה משפטים דלהלן:

\bigskip
\noindent\textbf{\underline{משפט 2.}}

\noindent כאשר
\[ m = ab \]

\noindent בתנאי
\[ (a,b) = 1 \]

\noindent אזי מספר השרשים היסודיים של הקונגרואנציה
\[ f(x) \equiv 0 \pmod{m} \]

\noindent שוה למכפלת מספר שרשי הקונגרואנציה
\[ f(x) \equiv 0 \pmod{a} \]

\noindent במספר שרשי הקונגרואנציה
\[ f(x) \equiv 0 \pmod{b} \, . \]

\noindent את ההוכחה למשפט זה אפשר למצוא בכל ספר לימוד בתורת המספרים
ועי"כ לא נביאה כאן.

% [original page: 48]
\bigskip
\noindent\textbf{\underline{משפט 3.}} לקונגרואנציה
\[ z^2 \equiv D \pmod{p^{\alpha}} \]

\noindent א) במקום ש-$p \neq 2$ הנו מספר ראשוני ו-$\alpha$ מספר טבעי כל שהוא ישנם
בדיוק שני שרשים כאשר
\[ p \nmid D \]

\noindent ו-$D$ הוא שארית ריבועית מודולו $p^{\alpha}$

\medskip
\noindent ב) במקום
\[ p = 2 \, , \quad \alpha = 1 \]

\noindent יש בדיוק שורש אחד ויחיד
\[ Z = 1 \]

\noindent כי במקרה זה
\[ 1 \equiv -1 \pmod{2} \]

\noindent כאשר
\[ 2 \nmid D \]

\noindent ויש שורש אחד ויחיד
\[ Z = 0 \]

\noindent כאשר
\[ 2 \mid D \]

\medskip
\noindent ג) במקום $p = 2 \, , \; \alpha > 2$ יש בדיוק \underline{ארבעה שרשים}, והם:
\[ \begin{array}{l}
Z_1 \equiv Z_0 \pmod{2^{\alpha}} \\
Z_2 \equiv -Z_0 \pmod{2^{\alpha}} \\
Z_3 \equiv Z_0 + 2^{\alpha-1} \pmod{2^{\alpha}} \\
Z_4 \equiv -Z_0 + 2^{\alpha-1} \pmod{2^{\alpha}}
\end{array} \]

\medskip
\noindent ד) במקום
\[ \alpha = 2 \, , \quad p = 2 \]

\noindent ישנם בדיוק \underline{שני שרשים}, והם:
\[ 1, \, 3 \, ; \qquad \text{או מה שהיינו הך:} \qquad 1, \, -1 \, , \]

\bigskip
\noindent\textbf{\underline{הערה 1.}} $Z_0$ הנו שורש יסודי כלשהו של הקונגרואנציה הנ"ל.

\noindent\textbf{\underline{הערה 2.}} בכל המקרים הנ"ל, אם $Z_0$ הנו שורש לקונגרואנציה הנ"ל,
אזי גם $-Z_0$ הנו שורש לאותה קונגרואנציה, פרט למקרה $m=2$ שבו
מתלכדים שני השרשים האלה, כי:
\[ 1 \equiv -1 \pmod{2} \]

% [original page: 49]
\noindent ככה, למשל, כאשר
\[ m = 8 \, , \quad D = 4k+1 \]

\noindent השרשים היסודיים של הקונגרואנציה
\[ z^2 \equiv D \pmod{m} \]

\noindent הם:
\[ 1, \, 3, \, 5, \, 7 \, , \]

\noindent וקיים:
\[ \left. \begin{array}{l}
7 \equiv -1 \pmod{8} \\
5 \equiv -3 \pmod{8}
\end{array} \right\} \]

\noindent דבר זה מתקיים גם בשרשי הקונגרואנציה
\[ z^2 \equiv D \pmod{2^{\alpha}} \]

\noindent ואמנם:
\[ -Z_0 + 2^{\alpha-1} = -Z_0 - 2^{\alpha-1} + 2^{\alpha} \]

\noindent ולכן:
\[ -Z_0 + 2^{\alpha-1} \equiv -Z_0 - 2^{\alpha-1} \pmod{2^{\alpha}} \]

\noindent כלומר, קיימים ארבעת השרשים:
\[ Z_0 \, , \; -Z_0 \]
\[ Z_0 + 2^{\alpha-1} \, , \; -Z_0 - 2^{\alpha-1} \]

\noindent\textbf{\underline{הערה 3.}} גם את ההוכחה למשפט האחרון לא נביא כאן.

\bigskip
\hrule
\bigskip

% [original page: 50]
\subsection*{2\S. פתרון הקונגרואנציה ממעלה שניה לפי גאוס:
ליקוי מתודי בשיטה זאת ותיקונה.}

\noindent 1. את הקונגרואנציה המלאה
\[ (1) \qquad ax^2 + bx + c \equiv 0 \pmod{m} \]

\noindent אנו פותרים, לפי גאוס, בדרך פשוטה פשוטה מאד.

\noindent מ-(1) נופק מיד כי:
\[ (2) \qquad (2ax+b)^2 \equiv D \pmod{4am} \]

\noindent ולכן קובע גאוס, עלינו לפתור את מערכת הקונגרואנציות
\[ (3) \qquad z^2 \equiv D \pmod{4am} \]
\[ 2ax + b \equiv Z \pmod{4am} \]

\noindent בכדי לקבל את השרשים היסודיים של הקונגרואנציה (1).

\noindent הקונגרואנציה הראשונה שייכת לתורת השאריות הריבועיות, שהז-
כרנו בסעיף הקודם, ואילו השנייה היא ממעלה ראשונה והתורה הזאת
ידועה היטב.

\noindent אולם, אם כי, שיטה זאת מועילה למציאת פתרון לקונגרואנציה
(1) הרי יש בה בכל זאת כמה ליקויים שיטתיים.

\noindent ליקויים אלה נובעים מהשימוש בטכניקה בתחביר מבין הקונגרואנציות (3)
בצורה שגאוס בחר, והדבר יתברר היטב בהמשך דיוננו על קונגרו-
אנציות.

\medskip
\noindent\textbf{\underline{2. קונגרואנציות דו-צדדיות.}}

\noindent בתורת התבניות אנו נתקלים בסוג מיוחד הנקרא בפי גאוס, דיריכלה
ואחרים בשם ``תבניות דו-צדדיות''.

\noindent והנה לפי הקלאספיקציה החדשה של קבוצות הפתרונות של תבניות אנו
מגיעים גם לקונגרואנציות מסוג השונה במקצת מן הרגילים, והמסקנות המת-
קבלות בסוג זה של קונגרואנציות שונות מאלה הקיימות בקונגרואנציות מהצורה
הרגילה. לפיכך נדון כאן בסוג זה.

% [original page: 51]
\bigskip
\noindent\textbf{\underline{הגדרה.}}

\noindent בשם קונגרואנציה דו-צדדית ממעלה שנייה נקרא לכל קונגרו-
אנציה
\[ (4) \qquad ax^2 + bx + c \equiv 0 \pmod{m} \]

\noindent אם קיים
\[ (5) \qquad a \mid b \, . \]

\noindent מכיוון ש
\[ (6) \qquad D = b^2 - 4ac \]

\noindent נופק מיד לפי (5), כי:
\[ (7) \qquad a \mid D \, ; \]

\noindent לפי הדרך של גאוס אנו מגיעים, איפוא, לפתרון מערכת הקונגרואנ-
ציות
\[ (8) \quad \left\{ \begin{array}{l}
z^2 \equiv D \pmod{4am} \\
2ax + b \equiv Z \pmod{4am}
\end{array} \right. \]

\noindent בתנאים שונים קצת מהכלליים, שפרטנו בסעיף הראשון. כי, לפי (7)
אין יותר מקום להנחה
\[ (D, m) = 1 \]

\noindent שומה עלינו, איפוא, לחקור עתה את המקרה האחרון.

\noindent נגביל את עצמנו לעת -- עתה להנחה
\[ (9) \qquad \boxed{\, a \text{ אינו מכיל גורמים ריבועיים} \,} \]

\noindent בתנאי זה נוכיח בשביל מטרותינו הבאות את המשפט הבא:

\bigskip
\noindent\textbf{\underline{משפט 4.}} לקונגרואנציה
\[ (10) \qquad z^2 \equiv D \pmod{4ap^{\alpha}} \]

\noindent ישנם בדיוק \underline{ארבעה שרשים יסודיים}, בתנאים הבאים:
\[ \begin{array}{ll}
\text{א)} & a \nmid D \, , \\
\text{ב)} & a = \text{לא זוגי} \\
\text{ג)} & p \neq 2 \text{ ראשוני} \\
\text{ד)} & D \equiv 0 \pmod{4} \\
\text{ה)} & a \text{ אינו מכיל גורמים ריבועיים} \\
\text{ו)} & p \nmid D \, .
\end{array} \]

% === TRANSCRIPTION NOTES ===
% General: the source typewriter consistently misprints "mod." as "mad." throughout
%   pages 44-51 (a typewriter/typist artifact, not a mathematical variable); all instances
%   were normalized to \pmod{...} per the style guide.
% Small pen/ink checkmarks appearing above certain letters (e.g. over "m_1" on pages 33-34,
%   and forming the Greek "nu" exponents "2ν-1" on pages 32,34-35) were treated as follows:
%   where the checkmark clearly stands in for the Greek letter nu (exponents like m^{2\nu-1},
%   m^{\nu-1}) it was transcribed as \nu; isolated checkmarks with no exponent role (pure
%   proofreading ticks) were treated as marginalia and ignored, per instructions.
% p032: "בעלת המבחין" and the phrase "במקום שהמספרים Z, n נקבעים" (page 32 / p040.png) --
%   the glyphs for מ/נ (m/n) and ה"טט"/"שפ" ("מטפט" vs "משפט") are visually similar on this
%   typewriter font; read as "משפט" (theorem) and "Z, n" respectively based on mathematical
%   context (n matches the third coefficient introduced later in the same derivation).
% p033-034 (p041-p042.png): several connective words in the statement/proof of Theorem 6
%   ("ולהיפך: כל זוג פתרונות ... שייכות לקבוצות הפתרונות המיוחסות") were reconstructed from
%   partially smudged type; the mathematical content (equations (1)-(8)) is legible and
%   transcribed with high confidence, but a few grammatical connector words are best-effort
%   readings rather than certain.
% p034 (p042.png): a stray "ה" catchword-like mark at the top right margin of the page
%   (above the continuation of equation from p041) was treated as a typographical/continuation
%   artifact and omitted, consistent with skipping marginalia.
% p036 (p044.png): the parenthetical "(בשביל כל m, ν)" is a best-effort reading of a
%   somewhat compressed inline phrase next to the boxed equation (16).
% p040 (p048.png): equation (18) reads "b_1 = -b - 2c\delta_1" with a raised checkmark over
%   the delta subscript numeral, ignored as marginalia (see general note above).
% p041 (p049.png): a small arrow-like pen mark appears after "b_\nu^2 < 4a_\nu^2" and again
%   after the boxed equation (23)-derivation on p042 (p050.png); these appear to be a
%   reader's/proofreader's flow-indicator marks rather than mathematical notation and were
%   omitted.
% p044 (p052.png): in the case-b derivation, the coefficient "2ax^{(\nu)}y^{(\nu)}" in the
%   source is printed without the \nu subscript on the first "a" (likely a typesetting typo
%   for $2a_\nu x^{(\nu)} y^{(\nu)}$); transcribed with the subscript restored for mathematical
%   consistency, following the pattern of the surrounding terms.
% p050 (p058.png): "וקבע גאוס" / "בטכניקה בתחביר" -- a short clause describing the
%   methodological defect in Gauss's approach to system (3) is faintly typed; transcribed
%   as best-effort ("ליקויים אלה נובעים מהשימוש בטכניקה בתחביר מבין הקונגרואנציות (3)
%   בצורה שגאוס בחר") -- the general sense (Gauss's chosen pairing/ordering of the two
%   congruences in (3) is methodologically flawed) is clear, but exact wording of this one
%   clause is not fully certain.
% No pages in this range (p040-p059) contained content judged genuinely illegible enough to
%   require an inline [?] marker; all pages were readable with the above noted caveats on a
%   handful of connective words/marginalia.
